% ============================================================
% Chapter 2: State of the Art
%   Section: Automated Fact-Checking
%     Subsection: The Fact-Checking Pipeline
% Included from chapters/02-state-of-the-art/01-automated-fact-checking/automated-fact-checking.tex
% ============================================================

\subsection{The Fact-Checking Pipeline}

Automated fact-checking is commonly framed as a pipeline that mirrors
the work of a human fact-checker \cite{thorne2018automated,
guo2022survey}. Guo, Schlichtkrull and Vlachos organise the field into
three stages and a fourth output:

\begin{itemize}
    \item \textbf{Claim detection}: deciding which sentences make a
    claim worth checking. ClaimBuster is the early reference: it scores
    sentences of political debates for \emph{check-worthiness}, the
    likelihood that a fact-checker would want to check them
    \cite{hassan2017claimbuster}.
    \item \textbf{Evidence retrieval}: finding the documents, and the
    passages in them, that bear on the claim.
    \item \textbf{Verdict prediction}: classifying the claim against
    the evidence, from a binary true/false to the fine-grained scales
    fact-checkers publish.
    \item \textbf{Justification}: producing an explanation of the
    verdict that a reader can check.
\end{itemize}

A recurring conclusion of the field is that full automation is not yet
realistic and that the most useful systems assist fact-checkers ---
finding claims, finding evidence, detecting claims already checked ---
rather than replace them \cite{nakov2021assisting}. This is the
question RQ4 puts to the system built here: not whether it replaces a
journalist, but which share of the work it can take on, at what
quality.

The proposed method follows the same decomposition
(Section~\ref{sec:design-fact-checking}). Claim selection plays the
part of claim detection, with a structural rather than a learned notion
of what is checkable; evidence retrieval and ranking find and filter
open-web documents; a language model predicts the verdict; and its
justification is the set of sources it cites, each with the verbatim
passage it relied on, which the system checks mechanically.
